\documentclass[11pt,a4paper]{article}

\usepackage[utf8]{inputenc}
\usepackage[vietnamese]{babel}
\usepackage{fontspec}
\setmainfont{Times New Roman}
\setsansfont{Arial}
\setmonofont{Consolas}

\usepackage[a4paper,top=1.8cm,bottom=2cm,left=1.8cm,right=1.8cm]{geometry}
\usepackage{amsmath,amssymb}
\usepackage{booktabs}
\usepackage{tabularx}
\usepackage{array}
\usepackage{listings}
\usepackage{xcolor}
\usepackage{fancyhdr}
\usepackage{multicol}
\usepackage{enumitem}
\usepackage{tcolorbox}
\usepackage{lastpage}

% Màu sắc
\definecolor{codegray}{rgb}{0.96,0.96,0.96}
\definecolor{keywordblue}{rgb}{0.0,0.2,0.6}
\definecolor{stringgreen}{rgb}{0.0,0.5,0.2}
\definecolor{commentgray}{rgb}{0.4,0.4,0.4}
\definecolor{correctgreen}{rgb}{0.0,0.5,0.15}
\definecolor{boxborder}{rgb}{0.7,0.7,0.7}

\lstset{
  language=Python,
  backgroundcolor=\color{codegray},
  basicstyle=\ttfamily\small,
  keywordstyle=\color{keywordblue}\bfseries,
  stringstyle=\color{stringgreen},
  commentstyle=\color{commentgray}\itshape,
  showstringspaces=false,
  breaklines=true,
  frame=single,
  rulecolor=\color{boxborder},
  aboveskip=4pt,
  belowskip=4pt,
  numbers=none,
  columns=fullflexible
}

% Header & Footer
\pagestyle{fancy}
\fancyhf{}
\setlength{\headheight}{15pt}
\lhead{\small \textbf{Học phần:} Phân tích dữ liệu với Python (DSAI1005)}
\rhead{\small \textbf{HƯỚNG DẪN CHẤM \& ĐÁP ÁN CHI TIẾT}}
\lfoot{\small Khoa Khoa học Dữ liệu \& Trí tuệ Nhân tạo -- ĐH Kinh tế Quốc dân}
\rfoot{\small Trang \thepage/\pageref{LastPage}}
\renewcommand{\headrulewidth}{0.5pt}
\renewcommand{\footrulewidth}{0.5pt}

\begin{document}

\noindent
\begin{minipage}[t]{0.55\textwidth}
  \centering
  \textbf{TRƯỜNG ĐẠI HỌC KINH TẾ QUỐC DÂN}\\
  \textbf{KHOA KHOA HỌC DỮ LIỆU \& AI}\\
  \rule{4cm}{0.4pt}\\[2mm]
  \textbf{HƯỚNG DẪN CHẤM \& ĐÁP ÁN CHÍNH THỨC}\\
  \textbf{Học phần:} Phân tích dữ liệu với Python (DSAI1005)\\
  \textbf{Nội dung:} Đề kiểm tra 60 phút (Tuần 1 -- 5) --- \textbf{Mã đề:} \textbf{101}
\end{minipage}
\hfill
\begin{minipage}[t]{0.42\textwidth}
  \begin{tcolorbox}[colback=blue!5!white,colframe=blue!75!black,title=\textbf{Tổng quan thang điểm 10}]
    \small
    \textbf{Phần I (Trắc nghiệm):} 20 câu $\times$ 0.25đ = \textbf{5.0 điểm}\\
    \textbf{Phần II (Tự luận ngắn):} 4 câu = \textbf{5.0 điểm}\\
    \rule{\linewidth}{0.3pt}\\
    \textbf{Tổng điểm toàn bài:} \textbf{10.0 điểm}
  \end{tcolorbox}
\end{minipage}

\vspace{3mm}
\hrule
\vspace{2mm}

\section*{PHẦN I: ĐÁP ÁN TRẮC NGHIỆM (5.0 điểm --- 0.25 điểm / câu)}

\begin{center}
\begin{tabular}{|c|c|c|c|c|c|c|c|c|c|c|}
  \hline
  \textbf{Câu} & \textbf{1} & \textbf{2} & \textbf{3} & \textbf{4} & \textbf{5} & \textbf{6} & \textbf{7} & \textbf{8} & \textbf{9} & \textbf{10} \\
  \hline
  \textbf{Đáp án} & \textbf{C} & \textbf{B} & \textbf{C} & \textbf{A} & \textbf{C} & \textbf{B} & \textbf{C} & \textbf{A} & \textbf{B} & \textbf{C} \\
  \hline
  \hline
  \textbf{Câu} & \textbf{11} & \textbf{12} & \textbf{13} & \textbf{14} & \textbf{15} & \textbf{16} & \textbf{17} & \textbf{18} & \textbf{19} & \textbf{20} \\
  \hline
  \textbf{Đáp án} & \textbf{A} & \textbf{A} & \textbf{B} & \textbf{B} & \textbf{B} & \textbf{A} & \textbf{C} & \textbf{A} & \textbf{B} & \textbf{A} \\
  \hline
\end{tabular}
\end{center}

\subsection*{Giải thích vắn tắt cơ sở đáp án trắc nghiệm:}
\begin{itemize}[leftmargin=*,itemsep=1.5pt]
  \item \textbf{Câu 1 (C):} Theo quy trình 6 bước chuẩn (Tuần 1): Xác định bài toán $\rightarrow$ Thu thập $\rightarrow$ Làm sạch $\rightarrow$ \textbf{EDA} $\rightarrow$ Mô hình hóa $\rightarrow$ Báo cáo/Trực quan hóa.
  \item \textbf{Câu 2 (B):} Dữ liệu lệch phải (right-skewed) do có outliers cực lớn thì Trung vị (Median) không bị ảnh hưởng, phản ánh trung thực hơn Trung bình (Mean).
  \item \textbf{Câu 3 (C):} $r = 0.88 > 0.7$ là tương quan tuyến tính thuận chiều rất mạnh (lưu ý tương quan không đồng nhất với quan hệ nhân quả tuyệt đối).
  \item \textbf{Câu 4 (A):} \texttt{a[1:, :2]}: dòng lấy từ index 1 đến hết (dòng 1, 2: giá trị 40, 50, 60 và 70, 80, 90); cột lấy vị trí 0, 1 $\rightarrow$ \texttt{[[40, 50], [70, 80]]}.
  \item \textbf{Câu 5 (C):} Quy tắc Broadcasting: \texttt{(3, 1)} ghép với \texttt{(1, 4)} mở rộng theo từng chiều thành \texttt{(3, 4)}.
  \item \textbf{Câu 6 (B):} \texttt{np.nanmean()} tự động loại bỏ các giá trị \texttt{NaN} trước khi chia trung bình: $(10 + 20 + 30)/3 = 20.0$.
  \item \textbf{Câu 7 (C):} Vectorization là cơ chế cốt lõi của NumPy chuyển phép toán trên mảng xuống thư viện C thực thi.
  \item \textbf{Câu 8 (A):} \texttt{.loc} dựa trên nhãn index (label-based), \texttt{.iloc} dựa trên vị trí số nguyên (integer position-based).
  \item \textbf{Câu 9 (B):} Trong Pandas, kết hợp nhiều điều kiện logic bắt buộc dùng toán tử bitwise \texttt{\&} và bọc từng điều kiện trong cặp ngoặc đơn \texttt{()}.
  \item \textbf{Câu 10 (C):} Nhóm \texttt{'South'} có 2 giá trị \texttt{Sales} là $200$ và $400$. Trung bình là $(200 + 400)/2 = 300.0$.
  \item \textbf{Câu 11 (A):} \texttt{df.pivot\_table()} xoay bảng từ dạng dài sang rộng và mặc định tính tổng hợp (aggregation).
  \item \textbf{Câu 12 (A):} Phép nối trái (\texttt{left join}) giữ nguyên toàn bộ các khóa từ bảng trái (\texttt{df1}); những khóa chỉ có ở bảng phải (\texttt{df2}) sẽ bị bỏ qua.
  \item \textbf{Câu 13 (B):} \texttt{Figure} là toàn bộ khung vẽ; còn \texttt{Axes} là vùng đồ thị con trực tiếp chứa dữ liệu vẽ, trục tọa độ và tiêu đề.
  \item \textbf{Câu 14 (B):} Trong Seaborn, tham số \texttt{hue} dùng để phân nhóm màu sắc theo một biến danh mục (categorical variable).
  \item \textbf{Câu 15 (B):} Box Plot (Biểu đồ hộp) thể hiện trực quan tứ phân vị $Q_1, Q_2, Q_3$ và các điểm ngoại lệ (outliers) nằm ngoài râu.
  \item \textbf{Câu 16 (A):} \texttt{sns.heatmap(corr, annot=True)} hiển thị ma trận nhiệt kèm nhãn số giá trị tương quan trong từng ô (\texttt{annot=True}).
  \item \textbf{Câu 17 (C):} \texttt{chunksize} trả về một iterator để đọc từng phần dữ liệu nhằm tối ưu bộ nhớ RAM.
  \item \textbf{Câu 18 (A):} Tham số \texttt{usecols} trong \texttt{pd.read\_csv()} nhận danh sách các cột cần nạp vào bộ nhớ.
  \item \textbf{Câu 19 (B):} \texttt{pd.json\_normalize()} là hàm chuyên dụng làm phẳng dữ liệu JSON phân cấp lồng nhau thành bảng 2D phẳng.
  \item \textbf{Câu 20 (A):} Cú pháp chuẩn kết nối và nạp SQL vào DataFrame là \texttt{pd.read\_sql\_query(sql, conn)}.
\end{itemize}

\vspace{3mm}
\hrule
\vspace{2mm}

\section*{PHẦN II: ĐÁP ÁN \& BIỂU ĐIỂM TỰ LUẬN NGẮN (5.0 điểm)}

\begin{tcolorbox}[colback=gray!5!white,colframe=gray!60!black,title=\textbf{Lưu ý chung dành cho Cán bộ chấm thi}]
\small
\begin{itemize}[leftmargin=*,itemsep=1pt]
  \item Không trừ điểm nếu sinh viên sử dụng dấu nháy kép \texttt{" "} thay cho dấu nháy đơn \texttt{' '} hoặc ngược lại.
  \item Không trừ điểm nếu sinh viên thiếu dấu đóng ngoặc tròn ở cuối dòng nhưng logic câu lệnh hoàn toàn chính xác.
  \item Các cách viết tương đương về mặt kết quả logic (ví dụ dùng cú pháp từ điển thay vì Named Aggregation trong Pandas) đều được hưởng trọn điểm.
\end{itemize}
\end{tcolorbox}

\subsection*{Câu 1: Thao tác Mảng NumPy (1.0 điểm)}
\begin{enumerate}[label=\textbf{\alph*)}]
  \item \textbf{(0.5 điểm):}
  \begin{itemize}
    \item \textbf{Tên cơ chế:} Cơ chế lan truyền / Mở rộng mảng (\textbf{Broadcasting}). \textit{(Đạt 0.25đ)}
    \item \textbf{Ma trận kết quả:} \textit{(Đạt 0.25đ)}
    $$\text{result} = \begin{bmatrix}
      11 & 27 & 43 \\
      16 & 32 & 48 \\
      21 & 37 & 53 \\
      6  & 17 & 28
    \end{bmatrix}$$
  \end{itemize}

  \item \textbf{(0.5 điểm):}
  \begin{itemize}
    \item Điều kiện \texttt{data[:, 0] >= 15} lọc ra các dòng có phần tử đầu tiên $\ge 15$, đó là dòng index 1 (giá trị đầu 15) và dòng index 2 (giá trị đầu 20).
    \item Lát cắt \texttt{1:} lấy từ cột index 1 đến hết (loại bỏ cột đầu).
    \item Kết quả in ra chính xác: \textit{(Đạt 0.5đ)}
    $$\begin{bmatrix}
      30 & 45 \\
      35 & 50
    \end{bmatrix}$$
  \end{itemize}
\end{enumerate}

\subsection*{Câu 2: Thao tác \& Gom nhóm Dữ liệu với Pandas (1.5 điểm)}
\begin{enumerate}[label=\textbf{\alph*)}]
  \item \textbf{(0.5 điểm):} Viết đúng dòng lệnh tạo cột mới:
\begin{lstlisting}
df_orders['net_price'] = df_orders['price'] * (1 - df_orders['discount'])
\end{lstlisting}

  \item \textbf{(1.0 điểm):} Gom nhóm và tổng hợp dữ liệu. Sinh viên có thể viết 1 trong 2 cách sau:
  \begin{itemize}
    \item \textit{Cách 1 (Khuyên dùng - Named Aggregation):}
\begin{lstlisting}
df_summary = df_orders.groupby('category').agg(
    total_revenue=('net_price', 'sum'),
    avg_discount=('discount', 'mean')
)
\end{lstlisting}
    \item \textit{Cách 2 (Dictionary Aggregation):}
\begin{lstlisting}
df_summary = df_orders.groupby('category').agg({
    'net_price': 'sum',
    'discount': 'mean'
}).rename(columns={'net_price': 'total_revenue', 'discount': 'avg_discount'})
\end{lstlisting}
    \item \textbf{Phân bổ điểm:} Đúng \texttt{groupby('category')} được 0.25đ; Đúng các hàm tổng hợp (\texttt{'sum'} và \texttt{'mean'}) được 0.5đ; Đổi đúng tên 2 cột kết quả được 0.25đ.
  \end{itemize}
\end{enumerate}

\subsection*{Câu 3: Truy xuất \& Lưu trữ Dữ liệu Đa nguồn (1.0 điểm)}
\begin{enumerate}[label=\textbf{\alph*)}]
  \item \textbf{(0.5 điểm):} Đọc file CSV với tham số:
\begin{lstlisting}
df_customers = pd.read_csv("data/customers.csv", parse_dates=['signup_date'], na_values=["?", "None", "NA"])
\end{lstlisting}
  \textit{Biểu điểm:} Đúng tham số \texttt{parse\_dates=['signup\_date']} được 0.25đ; Đúng tham số \texttt{na\_values=["?", "None", "NA"]} được 0.25đ.

  \item \textbf{(0.5 điểm):} Ghi DataFrame vào bảng SQLite:
\begin{lstlisting}
df_summary.to_sql("report_category", conn, if_exists="replace", index=False)
\end{lstlisting}
  \textit{Biểu điểm:} Đúng tên bảng và kết nối (\texttt{"report\_category", conn}) được 0.25đ; Đúng hai tham số (\texttt{if\_exists="replace", index=False}) được 0.25đ.
\end{enumerate}

\subsection*{Câu 4: Trực quan hóa Dữ liệu \& Rút ra Nhận định Phân tích (1.5 điểm)}
\begin{enumerate}[label=\textbf{\alph*)}]
  \item \textbf{(0.75 điểm):}
  \begin{itemize}
    \item \textbf{Loại biểu đồ:} Biểu đồ phân tán / Scatter plot. \textit{(Đạt 0.25đ)}
    \item \textbf{Mã lệnh Seaborn:} \textit{(Đạt 0.5đ)}
\begin{lstlisting}
sns.scatterplot(data=df, x='discount_pct', y='sales_volume', hue='customer_segment')
\end{lstlisting}
  \end{itemize}

  \item \textbf{(0.75 điểm) --- Nhận định kinh doanh:}
  \begin{itemize}
    \item \textit{Nhận định 1 (0.375đ):} Khách hàng phổ thông (\texttt{Standard}) rất nhạy cảm với giảm giá (độ co giãn cầu theo giá cao) --- mức chiết khấu càng cao thì lượng tiêu thụ tăng lên rõ rệt.
    \item \textit{Nhận định 2 (0.375đ):} Khách hàng \texttt{VIP} không nhạy cảm với giảm giá (mức mua ổn định bất kể khuyến mãi). Doanh nghiệp \textbf{không nên lãng phí ngân sách chiết khấu sâu cho nhóm VIP}, mà nên tập trung ưu đãi giảm giá để kích cầu nhóm Standard, đồng thời thiết kế chính sách chăm sóc/dịch vụ chuyên biệt (phi giá) cho nhóm VIP.
  \end{itemize}
\end{enumerate}

\begin{center}
\textbf{------------------- HẾT -------------------}
\end{center}

\end{document}
